\documentclass{beamer}
\usetheme{Madrid}

\title{Section 4.1: Vector Spaces}
\author{}
\date{}

\begin{document}

%------------------------------------------------
\begin{frame}
\titlepage
\end{frame}
%------------------------------------------------

\begin{frame}{What Is a Vector?}
Before this chapter:

\[
\begin{bmatrix}1 \\ 2 \\ 3\end{bmatrix}
\]

was a vector.

Now:

A \textbf{vector} is just an object we can:
\begin{itemize}
    \item Add
    \item Multiply by a number
\end{itemize}

If those operations behave nicely → it's a vector space.
\end{frame}

%------------------------------------------------

\begin{frame}{What Is a Vector Space?}
A vector space is a set $V$ where:

\begin{itemize}
    \item You can add vectors: $u + v$
    \item You can scale vectors: $c v$
    \item You stay inside the set (closure)
    \item There is a zero vector
    \item Every vector has a negative
\end{itemize}

If it behaves like $\mathbb{R}^n$, it's a vector space.
\end{frame}

%------------------------------------------------

\begin{frame}{Example 1: $\mathbb{R}^2$}
Vectors look like:

\[
\begin{bmatrix}x \\ y\end{bmatrix}
\]

Add:
\[
\begin{bmatrix}1 \\ 2\end{bmatrix}
+
\begin{bmatrix}3 \\ 4\end{bmatrix}
=
\begin{bmatrix}4 \\ 6\end{bmatrix}
\]

Scale:
\[
2 \begin{bmatrix}1 \\ 2\end{bmatrix}
=
\begin{bmatrix}2 \\ 4\end{bmatrix}
\]

Always stays in $\mathbb{R}^2$ → Vector space.
\end{frame}

%------------------------------------------------

\begin{frame}{Example 2: Polynomials}
Consider all polynomials of degree $\le 2$:

\[
a + bx + cx^2
\]

Example:
\[
1 + 2x + 3x^2
\]

Add two polynomials → still degree $\le 2$

Multiply by a number → still degree $\le 2$

So this is ALSO a vector space.
\end{frame}

%------------------------------------------------

\begin{frame}{Non-Example}
Set of vectors in $\mathbb{R}^2$ with only positive entries:

\[
\begin{bmatrix}1 \\ 1\end{bmatrix}
\]

But:

\[
-1 \cdot
\begin{bmatrix}1 \\ 1\end{bmatrix}
=
\begin{bmatrix}-1 \\ -1\end{bmatrix}
\]

Not in the set.

Fails closure → NOT a vector space.
\end{frame}

%------------------------------------------------

\begin{frame}{Why This Matters}
Now we can treat:

\begin{itemize}
    \item Vectors
    \item Polynomials
    \item Matrices
    \item Functions
\end{itemize}

all the same way.

All the ideas from Chapter 3 (span, independence, basis)
will work in ANY vector space.
\end{frame}

%------------------------------------------------

\begin{frame}{The Big Idea}
A vector space is:

\[
\text{A set where linear combinations make sense.}
\]

If you can do:

\[
c_1 v_1 + c_2 v_2 + \dots + c_k v_k
\]

and stay inside the set,

you are in a vector space.
\end{frame}

%------------------------------------------------

\end{document}
